

\subsection{The Case for Yearly Blind Test Sets}
\label{sec:manchester:qb-blind}

National Academic Quiz Tournaments (\abr{naqt}) professionalized these norms
and scaled them to large tournaments. Their goal is not only to write
questions, but to create efficient systems for determining which of many teams
knows the most. That requires both a discriminative format and a rigorous
answer-acceptability process.

Machine \abr{qa} has not yet matched this standard. Human competitions are
replenished continuously with new questions written and edited by domain
experts, while \abr{qa} leaderboards often reuse stale data written by
crowdworkers or harvested from the web. Even dynamic datasets such as
Dynabench, or our own TrickMe interface, are small compared to the tens of
thousands of questions written each year for human tournaments.

The natural remedy is a yearly blind test set. Systems should be evaluated once
on a fresh, hidden test set; then the questions and answers should be revealed,
and protests should be allowed. This would move \abr{qa} leaderboards closer to
the standards of human competitions and reduce the leakage problems discussed
earlier.



Taking our cue from the trivia community, we present an alternative
for \abr{mrqa}.
Blind test sets are created for a specific time; all systems are submitted simultaneously.
Then, all questions and answers are revealed.
System authors can protest correctness rulings on questions, directly
addressing the issues above. 
After agreement is reached,
quantitative metrics are computed for comparison purposes---despite their inherent limitations they at least can be
trusted.
Adopting this for \abr{mrqa} would require creating a new, smaller test
set every year.
However, this would gradually refine the annotations and process. xxzxzzzxzxzxxzxs 



\paragraph{How Adversarial Is Too Much?}
%
Common ground has its limits. That there is not a strict dichotomy between
Manchester and Cranfield can mask the importance of distinguishing motivations.
Some robustness proposals test models on inputs users would not make, or on
challenging but unnatural adversarial questions~\citep{Feng-18,jia2017adversarial,wallace-19,bartolo2020beat,kiela2021dynabench}.
These tests can be valuable for probing capabilities, but they do not directly
improve the user experience in the short term. The Manchester paradigm
welcomes these probes; the Cranfield paradigm will often ignore them. If we do
not state which goal we are pursuing, we risk building datasets that satisfy
neither.

\subsection{Focus on the bubble}

While every question should be perfect, time and resources are limited.  
Thus, authors and editors of tournaments ``focus on the bubble'', where the ``bubble'' are the questions most likely to discriminate between top teams at the tournament.
These questions are thoroughly playtested, vetted, and edited.
Only after these questions have been perfected will the other questions undergo the same level of polish.

For computers, the same logic applies.  
Authors should ensure that these discriminative questions are correct, free of ambiguity, and unimpeachable.
However, as far as we can tell, the authors of \qa{} datasets do not give any special attention to these questions.

Unlike a human trivia tournament, however---with finite patience of the participants---this does not mean that you should necessarily remove all of the easy or hard questions from your dataset.
This could inadvertently lead to systems unable to answer simple questions like ``who is buried in Grant's tomb?''~\cite[Chapter 7]{dwan-00}.
Instead, focus more resources on the bubble.

